\section{Admissibility Hierarchy on the $Q_8$ Cayley Graph}
  \label{sec:admissibility-hierarchy-on-the-q_8-cayley-graph}

  \begin{proposition}[Spectral admissibility hierarchy on $Q_8$]
    \label{prop:main}
    In the undirected Cayley graph $\Gamma(Q_8, \{\pm i, \pm j, \pm k\})$,
    under the local bounded-flux constraint~\eqref{eq:flux-bound} and linearised
    dispersion, the maximal admissible effective amplitude satisfies
    \begin{equation}
      A^{\max}_{\rho_4} = \frac{c_\chi}{\sqrt{6}}
      >
      \frac{c_\chi}{\sqrt{8}}
      = A^{\max}_{\rho_k},
      \qquad
      k = 1,2,3.
    \end{equation}
    The non-abelian sector $\rho_4$ occupies a strictly larger admissible region in the
    space of effective mode amplitudes than any nontrivial abelian sector, by a factor
    \begin{equation}
      \frac{A^{\max}_{\rho_4}}{A^{\max}_{\rho_k}}
      =
      \sqrt{\frac{8}{6}}
      =
      \sqrt{\frac{4}{3}}
      \approx 1.155.
    \end{equation}
  \end{proposition}

  \begin{proof}
    Direct substitution of $\lambda_{\rho_4} = 6$ and $\lambda_{\rho_k} = 8$
    into the admissibility envelope~\eqref{eq:envelope}.
  \end{proof}

  \begin{remark}
    The trivial sector $\rho_0$ has $\lambda = 0$ and thus formally unbounded $A^{\max}$.
    It corresponds to the uniform zero mode of the connected graph.
    In the present analysis it is factored out, since we are interested only in
    non-trivial excitations above the connected background.
  \end{remark}

  Proposition~\ref{prop:main}, together with
  Proposition~\ref{prop:basis-resolved}, establishes that, at fixed substrate capacity
  $c_\chi$, the non-abelian irreducible sector of the quaternion Cayley graph sustains
  a larger local effective amplitude before violating the bounded-flux constraint,
  even though its raw amplitude bound is lower than that of the nontrivial abelian
  sectors.

  In the language of the Cosmochrony programme:
  \begin{itemize}
    \item The admissibility envelope $A_n^{\max} = c_\chi / \sqrt{\lambda_n}$ defines a
    geometry of exclusion in the space of relational modes.
    Modes above the envelope are dynamically inadmissible.
    \item The non-abelian sector $\rho_4$ lies below this envelope for a larger
    range of effective amplitudes than the nontrivial abelian sectors
    $\rho_1, \rho_2, \rho_3$.
    The bounded-flux constraint is therefore less restrictive for non-abelian modes,
    enlarging their admissible window by a factor of $\sqrt{4/3}$.
    \item This advantage is a statement about the physically relevant local variable
    $A_n = a_n \|\psi_n\|_\infty$, not about the raw amplitude $a_n$.
    \item In the canonical basis, the weakly non-linear DBI correction again favours
    the non-abelian sector through its smaller eigenvalue.
    \item The exact one-mode DBI analysis shows that the same envelope
    $A_n^{\max} = c_\chi/\sqrt{\lambda_n}$ remains the true saturation boundary.
  \end{itemize}

  \begin{remark}[Real field and complex eigenmodes]
      The substrate field $\chi$ is taken real-valued.
      The natural eigenmodes of the 2D sector are matrix coefficients of $\rho_4$,
      which are complex-valued.
      One may either work in $\ell^2(Q_8,\mathbb{C})$ and take real and imaginary parts to
      obtain real eigenmodes, or directly choose a real basis for each eigenspace.
      The local admissibility criterion is intrinsically defined, but its expression in
      terms of $(a_n,\psi_n)$ depends on the chosen basis within a degenerate eigenspace
      through $\|\psi_n\|_\infty$.
      This basis dependence is resolved explicitly in
      Section~\ref{subsec:explicit-rho4}.
  \end{remark}
